\section{Conclusion and Future Work}

\begin{frame}{Conclusion}
    \begin{itemize}
        \item The proposal unifies emergence analysis, neuro-symbolic steering, and perturbation diagnostics in one local-interpretability program.
        \item The methodology jointly measures sparse feature geometry, cross-layer consistency, and logic satisfaction.
        \item Current evidence already supports feasibility for G1 and G3; G2 is positioned for final protocol-level validation.
    \end{itemize}
    \vspace{0.5em}
    \small
    Expected outcome: a practical, testable local analysis and intervention framework with explicit trade-offs and uncertainty reporting.
\end{frame}

\begin{frame}{Expected Contributions and Future Work}
    \small
    \begin{enumerate}
        \item Unified local-interpretability pipeline (SAE + $T_{\ell\to k}$ + LTN validators).
        \item Empirical evidence on emergence dynamics under controlled shifts.
        \item Measurable steering-locality framework with quality-vs-satisfaction curves.
        \item Transferable sparse diagnostics for perturbation/shift detection.
    \end{enumerate}
    \vspace{0.4em}
    \begin{block}{Immediate future work}
        Finalize frozen G2 runs, complete ablations and transfer tests, integrate results in final manuscript, and prepare defense-ready reproducibility artifacts.
    \end{block}
\end{frame}

\begin{frame}{Main Risks and Mitigations}
    \small
    \begin{table}
        \centering
        \begin{tabular}{@{}p{0.16\linewidth}p{0.34\linewidth}p{0.42\linewidth}@{}}
            \toprule
            \textbf{Risk} & \textbf{Impact} & \textbf{Mitigation} \\
            \midrule
            R1: Sparse feature instability &
            Weak reproducibility of local explanations &
            Multi-seed protocols, feature matching across runs, confidence intervals. \\
            \addlinespace[0.25em]
            R2: Weak semantic alignment across layers &
            Numeric fit without meaningful transfer &
            Validator-consistency checks and identity/baseline ablations. \\
            \addlinespace[0.25em]
            R3: Constraint gains with excessive quality loss &
            Steering not practical for deployment &
            Rule-weight calibration, Pareto reporting, unconstrained fallback baseline. \\
            \addlinespace[0.25em]
            R5: Time pressure near defense &
            Reduced ablation coverage &
            Freeze minimum publishable matrix and defer non-critical extensions. \\
            \bottomrule
        \end{tabular}
    \end{table}
\end{frame}
